\documentclass[10pt,a4paper,withhyper]{altacv}

\geometry{left=1.0cm,right=1.0cm,top=0.8cm,bottom=0.7cm,columnsep=0.8cm}

\usepackage{paracol}
\usepackage{enumitem}

% ---- Optimización de espacio (mismo estilo, ritmo vertical más compacto) ----
% Compactar listas
\setlist[itemize]{leftmargin=1.2em, topsep=1pt, partopsep=0pt, itemsep=1pt, parsep=0pt}

% Encabezado de sección con menos aire (bigskip -> medskip, medskip -> smallskip)
\renewcommand{\cvsection}[2][]{%
  \nointerlineskip\smallskip%
  \ifstrequal{#1}{}{}{\marginpar{\vspace*{\dimexpr1pt-\baselineskip}\raggedright\input{#1}}}%
  {\color{heading}\cvsectionfont\MakeUppercase{#2}}\\[-1.2ex]%
  {\color{headingrule}\rule{\linewidth}{2pt}\par}\smallskip
}

% Evento (puesto/proyecto) con menos espacio al cierre
\renewcommand{\cvevent}[4]{%
  {\large\color{emphasis}#1\par}
  \smallskip\normalsize
  \ifstrequal{#2}{}{}{%
  \textbf{\color{accent}#2}\par
  \smallskip}%
  \ifstrequal{#3}{}{}{%
    {\small\makebox[0.5\linewidth][l]%
      {\BeginAccSupp{method=pdfstringdef,ActualText={\datename:}}\cvDateMarker\EndAccSupp{}~#3}%
    }}%
  \ifstrequal{#4}{}{}{%
    {\small\makebox[0.5\linewidth][l]%
      {\BeginAccSupp{method=pdfstringdef,ActualText={\locationname:}}\cvLocationMarker\EndAccSupp{}~#4}%
    }}\par
  \smallskip\normalsize
}

% Separador entre proyectos más ajustado
\renewcommand{\divider}{\smallskip\textcolor{body!30}{\hdashrule{\linewidth}{0.6pt}{0.5ex}}\smallskip}

% Interlineado ligeramente más compacto
\linespread{0.95}

\iftutex
  \setmainfont{Roboto Slab}
  \setsansfont{Lato}
  \renewcommand{\familydefault}{\sfdefault}
\else
  \usepackage[rm]{roboto}
  \usepackage[defaultsans]{lato}
  \renewcommand{\familydefault}{\sfdefault}
\fi

\definecolor{SlateGrey}{HTML}{2E2E2E}
\definecolor{LightGrey}{HTML}{666666}
\definecolor{DarkBlue}{HTML}{0F3D5E}
\definecolor{AccentBlue}{HTML}{1D70A2}

\colorlet{name}{black}
\colorlet{tagline}{AccentBlue}
\colorlet{heading}{DarkBlue}
\colorlet{headingrule}{AccentBlue}
\colorlet{subheading}{AccentBlue}
\colorlet{accent}{AccentBlue}
\colorlet{emphasis}{SlateGrey}
\colorlet{body}{LightGrey}

\renewcommand{\namefont}{\Large\bfseries}
\renewcommand{\personalinfofont}{\footnotesize}
\renewcommand{\cvsectionfont}{\large\bfseries}

\setlength{\parskip}{0pt}
\setlength{\parindent}{0pt}

\begin{document}

\name{Lenin Adair Quezada Gómez}
\tagline{Ingeniería Financiera | Data Science | Quantitative Finance}

\personalinfo{
  \email{quezada.lenin@outlook.com}
  \phone{+52 33 14185861}
  \location{Guadalajara, Jalisco, México}
  \linkedin{lenin-quezada-6b77a63a8}
  \github{Lenon11}
}

\makecvheader

%==================== PERFIL (ancho completo) ====================

\cvsection{Perfil}

Estudiante de \textbf{Ingeniería Financiera en ITESO} con enfoque en \textbf{análisis financiero, riesgo crediticio y finanzas cuantitativas}, complementado con \textbf{ciencia de datos y machine learning}. Experiencia en el análisis de estados financieros de empresas públicas y en el desarrollo de modelos aplicados basados en datos. Interesado en roles donde pueda aplicar herramientas analíticas para resolver problemas reales.

\columnratio{0.72}

\begin{paracol}{2}

%==================== EXPERIENCIA ====================

\cvsection{Experiencia Laboral}

\cvevent{Analista Financiero}{World Vest Base}{Mayo 2026 -- Actualidad}{}
\begin{itemize}
\item Evaluación integral de empresas públicas, comprendiendo el análisis de sus estados financieros y de factores cualitativos determinantes para su desempeño.
\item Estructuración y depuración de información corporativa, garantizando la integridad y vigencia de la base de datos institucional.
\item Generación de reportes analíticos sobre la situación de cada empresa, proporcionando información estratégica que fundamenta las decisiones de inversión.
\end{itemize}

%==================== PROYECTOS ====================

\cvsection{Proyectos}

\cvevent{Mean-Reverting Portfolio}{PAP – Proyecto de Aplicación Profesional}{}{}
\begin{itemize}
\item Construcción de portafolios con grupos de activos cuyos precios tienden a moverse juntos y regresar a un valor de equilibrio.
\item Definición de reglas de compra y venta y validación de la estrategia con datos históricos en distintos periodos de mercado.
\item \textbf{Resultado:} alternativa de inversión de bajo riesgo con \textbf{$\sim$8.8\% de rendimiento anual}, Sharpe de \textbf{1.18} y máximo drawdown de 18\%.
\end{itemize}

\divider

\cvevent{Modelo de Riesgo Crediticio}{Proyecto académico y personal}{}{}
\begin{itemize}
\item Desarrollo de modelos para clasificar clientes según su riesgo de incumplimiento, con preparación de datos y selección de variables relevantes.
\item \textbf{Resultado:} score interpretable con \textbf{$\sim$78\% de accuracy} y \textbf{ROC-AUC de 0.75} en datos de prueba, bien calibrado frente a la tasa real de incumplimiento.
\end{itemize}

\divider

\cvevent{Valuación Crediticia de Empresas}{Proyecto académico presentado en Citi Banamex}{}{}
\begin{itemize}
\item Análisis de una empresa cotizada mediante indicadores financieros y comparación con empresas del mismo sector.
\item \textbf{Resultado:} recomendación de crédito con condiciones de financiamiento propuestas, presentada ante Citi Banamex.
\end{itemize}

\divider

\cvevent{Sistema de Visión Artificial}{}{}{}
\begin{itemize}
\item Desarrollo integral de un sistema de clasificación de objetos mediante \textbf{redes neuronales convolucionales (CNN)} (3 o más clases), con ingesta de un dataset propio de $\sim$4{,}000 imágenes y detección en tiempo real.
\item \textbf{Resultado:} precisión (accuracy) de aproximadamente \textbf{92\%}, con resultados sólidos y consistentes.
\end{itemize}

%==================== EDUCACIÓN ====================

\cvsection{Educación}

\cvevent{Ingeniería Financiera}{ITESO Universidad Jesuita de Guadalajara}{2022 -- Ago. 2026}{}
\begin{itemize}
\item \textbf{Promedio: 94.6 / 100.}
\item Áreas relevantes: finanzas cuantitativas, riesgo de crédito, econometría y análisis de datos.
\end{itemize}

%==================== COLUMNA DERECHA ====================

\switchcolumn

\cvsection{Habilidades}

\textbf{Lenguajes:} Python, SQL, MATLAB, TypeScript, HTML, CSS

\smallskip
\textbf{Herramientas:} React, Docker, Power BI, Git, GitHub, Linux, Terminal (CLI)

\smallskip
\textbf{Librerías:} Pandas, NumPy, Scikit-learn

\smallskip
\textbf{IA aplicada:} Claude Code, Codex e integración de IA en el flujo de trabajo

%==================== ÁREAS DE CONOCIMIENTO ====================

\cvsection{Áreas de Conocimiento}

\begin{itemize}
\item Análisis financiero
\item Riesgo crediticio
\item Análisis de datos
\item Finanzas cuantitativas
\item Machine learning aplicado
\end{itemize}

%==================== IDIOMAS ====================

\cvsection{Idiomas}

\textbf{Español:} Nativo

\smallskip
\textbf{Inglés:} Fluency English

\end{paracol}

\end{document}
